\documentclass[twocolumn]{aastex631}
\usepackage{amsmath}
\usepackage{booktabs}
\shorttitle{Selection-aware SDSS BPT/sSFR study}
\shortauthors{NebulaMind}
\begin{document}

\title{Broad Optical BPT Galaxies and Catalog Specific Star Formation in SDSS DR17: A Matched-Control Pilot}
\author{NebulaMind Research Autopilot}
\affiliation{Public SDSS DR17 data only}

\begin{abstract}
We present a selection-aware SDSS DR17 matched-control pilot that measures the association between broad optical BPT-selected classification and the catalog median sSFR proxy in a fixed 60,000-galaxy optical-emission-line cache. Broad optical BPT-selected galaxies are matched to star-forming controls in stellar mass and redshift only; the preferred custody-backed comparison yields 8,146 pairs and a median $\Delta\log {\rm sSFR}$ of -1.309 dex, with a bootstrap 95\% confidence interval of [-1.334,-1.283] dex. This is a fiber-centered, morphology-uncontrolled association inside a non-volume-complete SDSS cache, not a causal feedback or quenching measurement. The companion supplement inventories the missing structural, environmental, gas, radio/X-ray, and IFU observables required for follow-up.
\end{abstract}

\keywords{galaxies: active --- galaxies: star formation --- galaxies: evolution --- surveys --- methods: statistical}

\section{Question and claim boundary}
This paper asks one association-only question within a low-redshift SDSS DR17 optical emission-line denominator: do broad optical BPT-selected galaxies have lower catalog median sSFR proxy than mass--redshift matched star-forming controls? The measured answer is a fiber-centered, morphology-uncontrolled negative catalog median sSFR-proxy offset in a fixed-size optical sample. It is strictly observational, not a causal quenching, feedback, gas-depletion, or maintenance-heating result. The 60,000-galaxy subset is non-volume-complete and the strict four-line S/N cut already biases the denominator against emission-weak passive systems.

The scope excludes morphology or aperture controls, structural-proxy matching, Seyfert/LINER separation, bolometric accretion-luminosity proxies, gas-mass measurements, environment labels, and time-domain or duty-cycle modelling. BPT line ratios classify optical excitation, not black-hole accretion power; retired stellar populations and LIER-like emission can contaminate broad low-ionization classes and mimic active-nucleus signatures \citep{cidfernandes2011,stasinska2008,stasinska2015,belfiore2016}. Spatially resolved IFU data such as MaNGA or SAMI are therefore required to separate nuclear from extended low-ionization emission \citep{belfiore2016}. Because structural proxies were not retained in the cache, the present optical denominator cannot separate the measured offset from bulge-fraction, concentration-index, \texttt{fracDeV}, or central-velocity-dispersion associations. The analysis is restricted to $0.02<z<0.12$, so standard local BPT demarcations are used without redshift-evolution correction.

\subsection{Scope and limitations}
The association is defined inside a capped, selection-limited optical denominator. It is not a volume-complete census and does not include morphology, aperture fraction, group membership, halo mass, gas mass, or bolometric accretion-luminosity proxies as matching variables. The spectroscopy samples only the central 3-arcsec region (1.2--6.5 kpc here), so the offset remains degenerate with known correlations among stellar mass, bulge prominence, central velocity dispersion, and quenched fraction \citep{schawinski2010,bluck2014,piotrowska2022}. Single-fiber measurements can miss extended star-forming disks, so spatially resolved integral-field spectroscopy is required to break the aperture-morphology degeneracy \citep{penny2018,cheung2016,bundy2015,canodiaz2016}.

\section{Missing observables for future causal inference}
The present SDSS-only analysis is deliberately restricted to an optical association pilot, so the observables needed for any causal interpretation are not measured here. The remaining requirements are morphology and structural proxies, including Sersic index, bulge-to-total mass ratio, concentration, \texttt{fracDeV}, aperture-fraction control, group or halo membership, CO/HI gas masses, radio and X-ray proxies, resolved IFU kinematics, and matched simulation comparisons passed through the same selection function. These are the real-data inputs that the companion supplement inventories as future follow-up targets; they are not results of this paper.

\section{Data and shared selection}
The data backbone is public SDSS DR17 spectroscopy, photometry, emission-line measurements, and catalog physical-property estimates \citep{york2000,sdssdr17,brinchmann2004}. The retained pilot analysis sample is a fixed 60,000-galaxy subset selected sequentially by \texttt{specObjID}. It is a selection-limited pilot subset used to estimate the association within the available SDSS cache, not a volume-limited census. Because \texttt{specObjID} ordering follows SDSS targeting and plate/MJD bookkeeping, this non-random subset is not a random sky sample and introduces survey-plate and sky-coverage bias. The strict public four-line S/N$\geq3$ eligible parent count of 249,917 galaxies, and the corresponding 24.0\% cache coverage, are selection-context diagnostics rather than custody-backed independent result rows; the present custody receipt inventories the 60,000-row cache and matched-analysis artifacts, not the public SQL count logs behind the full parent cascade. Because the subset is fixed and non-volume-limited, it cannot be used to derive absolute volume densities, luminosity functions, or any population-normalized abundance. The redshift interval is not asserted to be uniformly mass-complete at its upper edge ($z=0.12$); completeness there remains a selection question rather than a retained result.
Over the redshift interval $0.02<z<0.12$, the SDSS 3-arcsec fiber subtends roughly 1.2--6.5 kpc, so the catalog median sSFR proxy comparison is fiber-centered rather than global.
Because the 3-arcsec fiber samples only the central regions at low redshift, the catalog-derived total sSFR proxy is an aperture-extrapolated quantity; the fixed 3-arcsec aperture systematically misses extended star-forming disks at low redshift \citep{kewley2005}. If broad optical BPT hosts are more bulge-dominated than the star-forming controls, the central fiber measurement can inflate the observed offset relative to a galaxy-wide star-formation comparison.
The stellar-mass and sSFR values are taken from the public MPA-JHU-style value-added table \texttt{galSpecExtra}, using its catalog median estimators \texttt{lgm\_tot\_p50} and \texttt{specsfr\_tot\_p50} after joining \texttt{SpecObj}, \texttt{galSpecInfo}, and \texttt{PhotoObj}. Although \texttt{PhotoObj} was joined in the catalog backbone, structural quantities such as \(R_{90}/R_{50}\), \texttt{fracDeV}, \texttt{petroR50}, and \texttt{petroR90} were not retained in the 60,000-galaxy cache, so morphology cannot be controlled in this cycle. Those are low-redshift SDSS catalog estimates, not rederived line-by-line physical measurements \citep{brinchmann2004,sdssdr17,york2000}. The catalog median sSFR proxy is therefore not interchangeable with a global UV+optical SFR estimate; comparisons to UV/optical SED-based SFRs are a required follow-up check for aperture-extrapolation systematics \citep{salim2007}. We use variance-normalized Euclidean matching, with each coordinate standardized by its sample standard deviation before distance calculation, because the feature space is only two variables, $(\log M_\star,z)$, so the rule stays transparent and the resulting nearest-neighbor control remains easy to interpret as an association baseline. A Mahalanobis distance would add an explicit covariance model for $(\log M_\star,z)$, but that extra layer is not warranted by this two-variable association pilot and would not change the boundary of the retained inference.

The selection is not neutral with respect to star formation. The strict four-line S/N cut preferentially removes emission-weak passive systems, while the sequential \texttt{specObjID} cap can inherit SDSS plate/MJD and sky-coverage structure. Public parent-count cascade values are therefore treated as query context only in this manuscript and are not promoted to retained results without separate query receipts. This caveat does not alter the inventoried 60,000-row cache or the 8,146-pair matched-offset result.

\section{Classification and matching}
BPT classes are computed from H$\alpha$, H$\beta$, [O~III]$\lambda5007$, and [N~II]$\lambda6584$ using standard demarcations \citep{baldwin1981,kewley2001,kauffmann2003bpt,kewley2006} (see Figure~\ref{fig:bpt}). The custody-backed analysis denominator contains 39,553 star-forming galaxies, 12,234 intermediate/composite galaxies, 8,146 broad optical BPT-selected targets, and 67 unclassified objects. The intermediate/composite galaxies are retained in denominator counts but are not part of the star-forming control pool used for matching. The 67 unclassified objects are retained in denominator counts for completeness but excluded from control pairing; they do not enter the 8,146-pair estimate because matching is performed only for broad optical BPT-selected targets with valid star-forming controls. Here, the star-forming control pool is defined as objects below the Kauffmann et al.\ (2003) demarcation, which is a conservative optical cut and may still include weak or obscured active-nucleus contaminants. Each broad optical BPT-selected galaxy is matched to the nearest star-forming control by variance-normalized Euclidean distance in standardized $(\log M_\star,z)$ space, with replacement. Because the two coordinates are standardized separately, the rule assigns equal variance weight to $\log M_\star$ and \(z\); this is transparent, but it can under- or over-penalize redshift separations relative to stellar-mass separations if the scientific cost of redshift mismatch differs from the sample variance. In the preferred estimate, this yields 100\% target coverage (8,146 of 8,146 targets matched), and the unrestricted Euclidean match has median absolute separations of 0.0045 dex in $\log M_\star$ and 0.00021 in redshift, so the association still inherits any mismatch in structure or fiber coverage between the two populations. No maximum mass--redshift caliper is imposed because the retained pilot estimate prioritizes complete coverage of all 8,146 custody-backed targets; imposing a caliper would create an additional selection function whose excluded targets are not separately inventoried in this candidate package. Matching is not performed in morphology, aperture fraction, halo mass, gas mass, accretion-luminosity proxy, or duty-cycle phase; these missing dimensions define follow-up requirements. A future real-data extension could use propensity-score or other multivariate balance diagnostics once those dimensions are actually available, but they are absent from this retained cache.
Here, ``broad optical BPT-selected'' means the inclusive optical-emission-line class under the standard BPT demarcations. Stricter Seyfert/LINER separation is a required follow-up analysis rather than a retained result in this provenance-limited manuscript; a future extension should apply explicit Seyfert/LINER demarcations such as the \citet{kewley2006} line-ratio boundaries and/or equivalent-width plus line-ratio classifications following \citet{cidfernandes2010}.

\begin{figure*}
\centering
\includegraphics[width=0.72\textwidth]{../figures/fig-bpt.pdf}
\caption{BPT line-ratio diagram for the SDSS DR17 analysis denominator. The matched controls are paired in stellar mass and redshift only, not in morphology, so the diagram verifies the optical-excitation classes used for matching but does not by itself establish an accretion-based interpretation.}
\label{fig:bpt}
\end{figure*}

\section{Matched-control result}
The preferred broad optical BPT comparison gives a large negative catalog median sSFR-proxy offset for the broad optical BPT-selected galaxies relative to star-forming controls. Without controlling for structural morphology or aperture fraction, a median $\Delta\log {\rm sSFR}$ (target minus matched control) of -1.309 dex is observed within this fixed-size, morphology-uncontrolled optical denominator and fiber-centered matched comparison. The robustness interval in Table~\ref{tab:robust} is a 95\% confidence interval on the median fiber-centered offset; the Scope and limitations subsection states the morphology and aperture limitations that govern its interpretation.

\begin{deluxetable*}{lrrrrrr}
\tabletypesize{\scriptsize}
\tablecaption{Provenance-retained matched catalog-sSFR offset.\label{tab:robust}}
\tablehead{\colhead{Target pool} & \colhead{Control pool} & \colhead{Match rule} & \colhead{$N$ pairs} & \colhead{Median $\Delta\log {\rm sSFR}$} & \colhead{95\% interval} & \colhead{Artifact IDs}}
\startdata
Broad optical BPT-selected targets & Star-forming controls below Kauffmann et al.\ (2003) & Variance-normalized Euclidean nearest neighbor in standardized $(\log M_\star,z)$ space, with replacement & 8,146 & -1.309 & [-1.334,-1.283] & \texttt{analysis\_results.json}; \texttt{matched\_agn\_sf\_pairs.csv} \\
\enddata
\tablecomments{$\Delta\log {\rm sSFR}$ is target minus matched star-forming control. This row is retained because it traces to \texttt{SDSS\_AGN\_SFR\_PILOT\_20260708T122000Z/analysis\_results.json} and \texttt{matched\_agn\_sf\_pairs.csv}, both inventoried in \texttt{provenance/REAL\_DATA\_SOURCE\_CUSTODY.json}. The preferred control rule is variance-normalized Euclidean nearest-neighbor matching in standardized $(\log M_\star,z)$ space, with replacement and no mass--redshift caliper. All values are conditional on the optical emission-line denominator.}
\end{deluxetable*}

\begin{figure*}
\centering
\includegraphics[width=0.86\textwidth]{../figures/fig-matched-offsets.pdf}
\caption{Distribution of matched-pair catalog-sSFR offsets for the preferred broad optical BPT-selected galaxy minus nearest star-forming control estimate ($N=8{,}146$ pairs, variance-normalized Euclidean matching in standardized $(\log M_\star,z)$ space, with replacement and without a maximum mass--redshift caliper). The median offset is $-1.309$ dex with bootstrap 95\% interval $[-1.334,-1.283]$ dex. The present custody file does not inventory separate artifacts for caliper, no-replacement, stricter-S/N, or Seyfert-like sensitivity variants.}
\label{fig:offsets}
\end{figure*}

\section{Interpretation}
The result is directly measured in the capped sample and remains falsifiable within the stated denominator. Figure~\ref{fig:offsets} shows the matched-offset distribution. Because the comparison is fiber-centered and selection-limited, the offset is a denominator-level association statement rather than a galaxy-wide causal inference. Any mechanistic interpretation still requires morphology, aperture controls, Seyfert/LINER separation, bolometric accretion-luminosity proxy, gas mass, environment, and time-domain/duty-cycle modelling. Spatially resolved IFU data are also needed to distinguish nuclear low-ionization excitation from extended LIER-like emission before the broad BPT class is interpreted as an accretion-selected LINER/Seyfert population \citep{belfiore2016}. Even central-velocity-dispersion control would not by itself establish causality, because fiber-centered SFR proxies can retain systematic offsets relative to broader star-formation estimates \citep{salim2012}.

\section{Conclusion}
RP-1 is a selection-aware pilot association paper. The fixed 60,000-galaxy subset is sequentially selected by \texttt{specObjID} and lacks morphological, structural, and aperture-fraction controls. The retained result is the 8,146-pair, -1.309 dex offset with bootstrap 95\% interval [-1.334,-1.283] dex. The accompanying \emph{Supplementary SDSS Denominator and Proxy Atlas for Galaxy-Evolution Follow-up} lists the missing observables required for future real-data tests. No measured result in this paper should be read as a gas-mass, environment-density, or feedback-efficiency estimate.
Future physical validation requires the kinds of measurements used in radio-mode and X-ray maintenance-heating studies \citep{best2005,fabian2012,mcnamara2007,heckmanbest2014,lamassa2013}, molecular and neutral gas studies \citep{xcoldgass2017,xgass2018}, outflow and kinematic studies \citep{veilleux2005,cicone2014,carniani2017,fiore2017}, and simulation comparisons passed through the same selection functions \citep{simba2019,tng2019,eagle2015}, together with the environment/context references \citep{peng2010,ellison2011,piotrowska2022,wetzel2013,dekel2006}. These references are examples of missing observables for future follow-up, not validation of any mechanism in this SDSS-only denominator. The reported offset remains association-only until morphology, aperture fraction, and the missing multiwavelength or IFU observables are added.

\section*{Data Availability}
This paper uses public SDSS DR17 spectroscopy, photometry, emission-line measurements, and MPA-JHU-style value-added catalog tables only. No proprietary data were used. The custody receipt is \texttt{provenance/REAL\_DATA\_SOURCE\_CUSTODY.json}. It inventories the retained source artifacts without copying the source data, including \texttt{SDSS\_AGN\_SFR\_PILOT\_20260708T122000Z/analysis\_results.json} (SHA-256 \texttt{668ad7a67290600ff5028ae587d32ef239a09bd8627a480539f37e1927d659df}), \texttt{data/analysis\_sample\_bpt.csv} (60,000 rows; SHA-256 \texttt{6f982fa5778c3900239149b28729f701390fe393a164b95236229adc1e422883}), and \texttt{data/matched\_agn\_sf\_pairs.csv} (8,146 rows; SHA-256 \texttt{4ea53af867cccccb2b68b81557ff84fe90ec3f13e0512ffbdc977fa7216996fd}). These artifacts support the headline 8,146-pair, -1.309 dex association estimate and its JSON-stored bootstrap interval. This candidate package does not contain the executable analysis script or frozen command recipe that generated the downstream JSON and CSV files, so full regeneration from the candidate alone is not claimed. No mock, synthetic, fake, placeholder, or toy data were used.

\facilities{SDSS}

\begin{thebibliography}{}
\bibitem[Abdurro'uf et al.(2022)]{sdssdr17} Abdurro'uf, Accetta, K., Aerts, C., et al. 2022, ApJS, 259, 35; ADS bibcode: 2022ApJS..259...35A; DOI: 10.3847/1538-4365/ac4a0c
\bibitem[Baldwin et al.(1981)]{baldwin1981} Baldwin, J.~A., Phillips, M.~M., \& Terlevich, R. 1981, PASP, 93, 5; ADS bibcode: 1981PASP...93....5B; DOI: 10.1086/130766
\bibitem[Best et al.(2005)]{best2005} Best, P.~N., Kauffmann, G., Heckman, T.~M., et al. 2005, MNRAS, 362, 25; ADS bibcode: 2005MNRAS.362...25B; DOI: 10.1111/j.1365-2966.2005.09192.x
\bibitem[Belfiore et al.(2016)]{belfiore2016} Belfiore, A., Maiolino, R., Maraston, C., et al. 2016, MNRAS, 461, 3111; ADS bibcode: 2016MNRAS.461.3111B; DOI: 10.1093/mnras/stw1234
\bibitem[Bluck et al.(2014)]{bluck2014} Bluck, A.~F.~L., Bruce, V.~A., Pilkington, K., et al. 2014, MNRAS, 441, 599; ADS bibcode: 2014MNRAS.441..599B; DOI: 10.1093/mnras/stu504
\bibitem[Brinchmann et al.(2004)]{brinchmann2004} Brinchmann, J., Charlot, S., White, S.~D.~M., et al. 2004, MNRAS, 351, 1151; ADS bibcode: 2004MNRAS.351.1151B; DOI: 10.1111/j.1365-2966.2004.08017.x
\bibitem[Bundy et al.(2015)]{bundy2015} Bundy, K., Law, D.~R., Yan, R., et al. 2015, ApJ, 798, 7; ADS bibcode: 2015ApJ...798....7B
\bibitem[Cano-D\'{\i}az et al.(2016)]{canodiaz2016} Cano-D\'{\i}az, M., Maiolino, R., Marconi, A., et al. 2016, ApJ, 821, L26; ADS bibcode: 2016ApJ...821L..26C; DOI: 10.3847/2041-8205/821/2/L26
\bibitem[Cid Fernandes et al.(2011)]{cidfernandes2011} Cid Fernandes, R., Stasi{\'n}ska, G., Schlickmann, S., et al. 2011, MNRAS, 413, 1687; ADS bibcode: 2011MNRAS.413.1687C; DOI: 10.1111/j.1365-2966.2011.18244.x
\bibitem[Cheung et al.(2016)]{cheung2016} Cheung, E., Bundy, K., Cappellari, M., et al. 2016, Nature, 533, 504; ADS bibcode: 2016Natur.533..504C; DOI: 10.1038/nature18006
\bibitem[Cid Fernandes et al.(2010)]{cidfernandes2010} Cid Fernandes, R., Stasi{\'n}ska, G., Schlickmann, M.~S., Mateus, A., Vale Asari, N., Schoenell, W., \& Sodr{\'e}, L. 2010, MNRAS, 403, 1036; ADS bibcode: 2010MNRAS.403.1036C; DOI: 10.1111/j.1365-2966.2010.16486.x
\bibitem[Ellison et al.(2011)]{ellison2011} Ellison, S.~L., Patton, D.~R., Mendel, J.~T., et al. 2011, MNRAS, 418, 2043; ADS bibcode: 2011MNRAS.418.2043E; DOI: 10.1111/j.1365-2966.2011.19624.x
\bibitem[Ellison et al.(2021)]{ellison2021} Ellison, S.~L., Lin, L., Rosario, D.~J., et al. 2021, MNRAS, 501, 4777; ADS bibcode: 2021MNRAS.501.4777E; DOI: 10.1093/mnras/staa3846
\bibitem[Harrison(2017)]{harrison2017} Harrison, C.~M. 2017, Nature Astronomy, 1, 0165; ADS bibcode: 2017NatAs...1E.165H; arXiv: 1703.06889
\bibitem[Carniani et al.(2017)]{carniani2017} Carniani, S., Marconi, A., Maiolino, R., et al. 2017, A\&A, 605, A42; ADS bibcode: 2017A\&A...605A..42C; DOI: 10.1051/0004-6361/201730937
\bibitem[Catinella et al.(2018)]{xgass2018} Catinella, B., Saintonge, A., Janowiecki, S., et al. 2018, MNRAS, 476, 875; ADS bibcode: 2018MNRAS.476..875C
\bibitem[Cicone et al.(2014)]{cicone2014} Cicone, C., Maiolino, R., Sturm, E., et al. 2014, A\&A, 562, A21; ADS bibcode: 2014A\&A...562A..21C; DOI: 10.1051/0004-6361/201322464
\bibitem[Dav{\'e} et al.(2019)]{simba2019} Dav{\'e}, R., Angl{\'e}s-Alc{\'a}zar, D., Narayanan, D., et al. 2019, MNRAS, 486, 2827; ADS bibcode: 2019MNRAS.486.2827D; DOI: 10.1093/mnras/stz937
\bibitem[Dekel \& Birnboim(2006)]{dekel2006} Dekel, A., \& Birnboim, Y. 2006, MNRAS, 368, 2; ADS bibcode: 2006MNRAS.368....2D; DOI: 10.1111/j.1365-2966.2006.10145.x
\bibitem[Fabian(2012)]{fabian2012} Fabian, A.~C. 2012, ARA\&A, 50, 455; ADS bibcode: 2012ARA\&A..50..455F; DOI: 10.1146/annurev-astro-081811-125521
\bibitem[Fiore et al.(2017)]{fiore2017} Fiore, F., Feruglio, C., Shankar, F., et al. 2017, A\&A, 601, A143; ADS bibcode: 2017A\&A...601A.143F; DOI: 10.1051/0004-6361/201629478
\bibitem[Heckman \& Best(2014)]{heckmanbest2014} Heckman, T.~M., \& Best, P.~N. 2014, ARA\&A, 52, 589; ADS bibcode: 2014ARA\&A..52..589H; DOI: 10.1146/annurev-astro-081913-035933
\bibitem[Kauffmann et al.(2003a)]{kauffmann2003bpt} Kauffmann, G., Heckman, T.~M., Tremonti, C., et al. 2003a, MNRAS, 346, 1055; ADS bibcode: 2003MNRAS.346.1055K
\bibitem[Kewley et al.(2001)]{kewley2001} Kewley, L.~J., Dopita, M.~A., Sutherland, R.~S., Heisler, C.~A., \& Trevena, J. 2001, ApJ, 556, 121; ADS bibcode: 2001ApJ...556..121K; DOI: 10.1086/321545
\bibitem[Kewley et al.(2005)]{kewley2005} Kewley, L.~J., Jansen, R.~A., \& Geller, M.~J. 2005, PASP, 117, 227; ADS bibcode: 2005PASP..117..227K; DOI: 10.1086/428303
\bibitem[Kewley et al.(2006)]{kewley2006} Kewley, L.~J., Groves, B., Kauffmann, G., \& Heckman, T. 2006, MNRAS, 372, 961; ADS bibcode: 2006MNRAS.372..961K; DOI: 10.1111/j.1365-2966.2006.10859.x
\bibitem[LaMassa et al.(2013)]{lamassa2013} LaMassa, S.~M., Heckman, T.~M., Ptak, A., \& Urry, C.~M. 2013, ApJL, 765, L33; ADS bibcode: 2013ApJ...765L..33L; DOI: 10.1088/2041-8205/765/2/L33
\bibitem[McNamara \& Nulsen(2007)]{mcnamara2007} McNamara, B.~R., \& Nulsen, P.~E.~J. 2007, ARA\&A, 45, 117; ADS bibcode: 2007ARA\&A..45..117M; DOI: 10.1146/annurev.astro.45.051806.110625
\bibitem[Nelson et al.(2019)]{tng2019} Nelson, D., Springel, V., Pillepich, A., et al. 2019, Computational Astrophysics and Cosmology, 6, 2; ADS bibcode: 2019CoAst...6....2N; DOI: 10.1186/s40668-019-0028-x
\bibitem[Penny et al.(2018)]{penny2018} Penny, S.~J., Davies, R.~L., Houghton, R.~C.~W., et al. 2018, MNRAS, 476, 979; ADS bibcode: 2018MNRAS.476..979P; DOI: 10.1093/mnras/sty222
\bibitem[Peng et al.(2010)]{peng2010} Peng, Y.-j., Lilly, S.~J., Kovac, K., et al. 2010, ApJ, 721, 193; ADS bibcode: 2010ApJ...721..193P; DOI: 10.1088/0004-637X/721/1/193
\bibitem[Piotrowska et al.(2022)]{piotrowska2022} Piotrowska, J.~M., Bluck, A.~F.~L., Maiolino, R., \& Peng, Y.-j. 2022, MNRAS, 512, 1052; ADS bibcode: 2022MNRAS.512.1052P; DOI: 10.1093/mnras/stac532
\bibitem[Schawinski et al.(2010)]{schawinski2010} Schawinski, K., Evans, D.~A., Virani, S., et al. 2010, ApJ, 711, 284; ADS bibcode: 2010ApJ...711..284S; DOI: 10.1088/0004-637X/711/1/284
\bibitem[Saintonge et al.(2017)]{xcoldgass2017} Saintonge, A., Catinella, B., Tacconi, L.~J., et al. 2017, ApJS, 233, 22; ADS bibcode: 2017ApJS..233...22S
\bibitem[Schaye et al.(2015)]{eagle2015} Schaye, J., Crain, R.~A., Bower, R.~G., et al. 2015, MNRAS, 446, 521; ADS bibcode: 2015MNRAS.446..521S; DOI: 10.1093/mnras/stu2058
\bibitem[Strateva et al.(2001)]{strateva2001} Strateva, I., Ivezi{\'c}, {\v Z}., Knapp, G.~R., et al. 2001, AJ, 122, 1861; ADS bibcode: 2001AJ....122.1861S; DOI: 10.1086/323301
\bibitem[Mendel et al.(2014)]{mendel2014} Mendel, J.~T., Simard, L., Palmer, M., Ellison, S.~L., \& Patton, D.~R. 2014, ApJS, 210, 3; ADS bibcode: 2014ApJS..210....3M; DOI: 10.1088/0067-0049/210/1/3
\bibitem[Salim et al.(2007)]{salim2007} Salim, S., Rich, R.~M., Charlot, S., et al. 2007, ApJS, 173, 267; DOI: 10.1086/519218
\bibitem[Salim et al.(2012)]{salim2012} Salim, S., Fang, J.~J., Rich, R.~M., et al. 2012, ApJ, 755, 105; DOI: 10.1088/0004-637X/755/2/105
\bibitem[Stasi{\'n}ska et al.(2008)]{stasinska2008} Stasi{\'n}ska, G., Asari, N.~V., Cid Fernandes, R., et al. 2008, MNRAS, 391, L29; ADS bibcode: 2008MNRAS.391L..29S; DOI: 10.1111/j.1745-3933.2008.00550.x
\bibitem[Stasi{\'n}ska et al.(2015)]{stasinska2015} Stasi{\'n}ska, G., Costa Duarte, M.~V., Vale Asari, N., Cid Fernandes, R., \& Sodr{\'e}, L. 2015, MNRAS, 449, 559; ADS bibcode: 2015MNRAS.449..559S; DOI: 10.1093/mnras/stv353
\bibitem[Veilleux et al.(2005)]{veilleux2005} Veilleux, S., Cecil, G., \& Bland-Hawthorn, J. 2005, ARA\&A, 43, 769; ADS bibcode: 2005ARA\&A..43..769V; DOI: 10.1146/annurev.astro.43.072103.150242
\bibitem[Wetzel et al.(2013)]{wetzel2013} Wetzel, A.~R., Tinker, J.~L., Conroy, C., \& van den Bosch, F.~C. 2013, MNRAS, 432, 336; ADS bibcode: 2013MNRAS.432..336W; DOI: 10.1093/mnras/stt410
\bibitem[York et al.(2000)]{york2000} York, D.~G., Adelman, J., Anderson, J.~E., Jr., et al. 2000, AJ, 120, 1579; ADS bibcode: 2000AJ....120.1579Y; DOI: 10.1086/301513
\end{thebibliography}

\end{document}
